% figuras/tikz/cap01-linea-marketing-x0.tex
% Línea de tiempo crítica: la serie "Marketing X.0" de Kotler, Kartajaya y Setiawan
% como periodización editorial (no modelo académico validado empíricamente).
\begin{tikzpicture}[
    hito/.style={rectangle, rounded corners, draw=guindaIPN, thick, fill=guindaIPNClara,
                 minimum width=2.6cm, minimum height=1.4cm, align=center, font=\footnotesize},
    linea/.style={thick, draw=grisIPN}
  ]
  \draw[linea, -{Stealth[length=3mm]}] (-0.3,0) -- (13.3,0);

  \node[hito] (m1) at (0,0.9) {\textbf{1.0}\\[1pt] Producto\\ (s.\,f.)};
  \node[hito] (m2) at (3.3,0.9) {\textbf{2.0}\\[1pt] Consumidor\\ (s.\,f.)};
  \node[hito] (m3) at (6.6,0.9) {\textbf{3.0}\\[1pt] Valores\\ 2010};
  \node[hito] (m4) at (9.9,0.9) {\textbf{4.0}\\[1pt] Digital\\ 2016};
  \node[hito] (m5) at (13.2,0.9) {\textbf{5.0}\\[1pt] Tecnología\\ 2021};

  \foreach \x in {0,3.3,6.6,9.9,13.2} {
    \draw[linea] (\x,0) -- (\x,0.2);
  }

  \node[hito, fill=bronceIPNClaro, draw=bronceIPN] (m6) at (13.2,-1.3) {\textbf{6.0}\\[1pt] Inmersivo\\ 2024};
  \draw[linea] (13.2,-0.2) -- (13.2,-0.6);

  \node[align=center, font=\scriptsize\itshape, text width=13.5cm] at (6.5,-2.6)
    {Serie editorial de libros (Kotler, Kartajaya y Setiawan), no un programa de\\
     investigación arbitrado: cada entrega se alinea a la tecnología dominante de su año de publicación.};
\end{tikzpicture}
